\clearpage\begingroup\SourceLeaf{22}\HistoricalHead{17}{R}{GAUSS AN MÖBIUS.}
\begin{center}
[2.]\\[.7em]
\textsc{Gauss} an \textsc{Möbius}.
\par\medskip\rule{\textwidth}{.35pt}\par
{\small [Berichte der mathematisch\,-\,physischen Classe der Königl. Sächsischen Gesellschaft der Wissenschaften zu Leipzig \textbf{31}, 1879, S. 61—64.]}\par
\rule{\textwidth}{.35pt}
\end{center}
\noindent\hspace*{.2\textwidth}Hochgeschätztester Freund!
\par\noindent\hspace*{\parindent}Für Ihr so schätzbares Werk über die Statik[\textsuperscript{*})]\begingroup\renewcommand{\thefootnote}{}\footnotetext{[\textsuperscript{*}) A. F. \textsc{Möbius}, \textit{Lehrbuch der Statik}, 2 Bände, Leipzig 1837, Gesammelte Werke III, 1886.]}\endgroup, sowie für Ihr gütiges
Schreiben sage ich Ihnen meinen verbindlichsten Dank. Unter den ununterbrochenen
Zerstreuungen unserer Festwoche habe ich gar keine, und in den
nächstfolgenden Tagen habe ich nur erst wenige Stunden frei gehabt, und
Ihr Werk eigentlich nur erst ein wenig durchblättern können. Indessen genug,
um Ihre gewohnte Gründlichkeit und Gediegenheit zu erkennen; mich
verlangt nach der Zeit, wo ich es mit Musse werde lesen können.

Anmerkungen \textls[120]{über Ihr Werk} kann ich Ihnen demnach jetzt noch gar
keine geben. Ich will jedoch, Ihnen mein Interesse an dem Gegenstande
selbst zu zeigen, eine kleine Bemerkung über das Princip der virtuellen Geschwindigkeiten,
oder wie ich lieber spreche, das Princip der virtuellen Bewegungen
beifügen, welcher Sie aber das Flüchtige und Ungeordnete der
Einkleidung nachsehen müssen.

Ich kann nicht leugnen, dass ich mit der Art, wie dieses Princip gewöhnlich
dargestellt wird, nicht ganz zufrieden bin; selbst der Begriff der
virtuellen Geschwindigkeiten oder besser der virtuellen Bewegungen, scheint
mir nicht alle zu wünschende Präcision und Rundung zu haben.

Ich habe ein Paarmahl Veranlassung genommen, das herrliche Princip
in derjenigen Gestalt, wie ich es wohl bei mündlichen Vorträgen darzustellen
pflege, auszusprechen, z. B. im ersten Artikel meiner \textit{Theoria generalis figurae
fluidorum}[\textsuperscript{**})]\begingroup\renewcommand{\thefootnote}{}\footnotetext{[\textsuperscript{**}) Werke V, S. 35.]}\endgroup. Es scheint dies aber wenig beachtet zu sein, vermuthlich, weil
es nur so beiläufig vorkommt, ohne das Unterscheidende besonders hervorzuheben
oder umständlich zu motiviren. Nur einmal erinnere ich, mit Herrn
Prof. \textsc{Dirksen} aus Berlin über \textls[120]{meine} Fassung gesprochen zu haben, wobei
\HistoricalFoot{XI{\scriptsize 1}.}{3}
\clearpage\SourceLeaf{23}\HistoricalHead{18}{L}{PHYSIK. BRIEFWECHSEL.}
\noindent ich die Satisfaction hatte ihn zu überzeugen, was freilich durch eine mündliche
Discussion leichter geschieht, als durch ein Paar in der Eile geschriebene
Zeilen.

Der Name \textls[120]{virtuelle} Bewegung scheint mir etwas ungeschickt gewählt
zu sein; indessen ist er einmal recipirt. Am a. O. finden Sie meine scharfe
Bestimmung des Begriffs. Wäre der Name nicht einmal da, so wäre vielleicht
\textls[120]{facultative} Bewegungen ein schicklicherer, und das \textls[120]{worauf es eigentlich
ankommt}, bestimmter bezeichnender. Doch in \textls[120]{verbis simus faciles}.

Wichtiger aber erscheint mir, dass alle Schriftsteller, auch \textsc{Lagrange}, das
Princip enger gefasst haben als nöthig ist. Allen Schriftstellern sieht man
die Fessel anhängen, die sie in den engern Raum festbannt, wo nur von
solchen Fällen die Rede ist, die Einer analytischen Formel vollständig unterworfen
sind.

Dies aber ist nicht die Erschöpfung der Fälle der \textls[120]{Natur}. Diese bietet
ebenso oft Fälle dar, wo die Eine analytische Formel [nicht] ausreicht. Ist
es dann nicht ein grosser Gewinn, wenn man das höchste Princip der Statik
\textls[120]{so} ausspricht, dass es die wirklichen Fälle der Natur in Einem Schlage unter
sich hat?

Ich erkläre mich etwas deutlicher.

Unter den Bedingungen, denen die facultative Bewegung eines Systems
materieller Punkte unterworfen ist, führt man z. B. wohl an, dass einer der
materiellen Punkte sich nur in einer bestimmten Fläche bewegen kann, setzt
dies in Gleichung u.s.w. Allein in der Natur ist dies fast nie so; vielmehr
gewöhnlich so, dass die Fläche \textls[120]{quaest.} nur die Oberfläche eines materiellen
Solidum ist, dessen Impenetrabilität den materiellen Punkt \textls[120]{nur} hindert, \textls[120]{in}
das Solidum einzudringen, aber nicht, sich davon zu entfernen.

Ein zweites Beispiel wäre, dass zwei materielle Punkte eine unveränderliche
Distanz von einander behalten müssen. Aber in der Natur ist es
äusserst häufig nur so (wenn z. B. die Verbindung durch einen biegsamen
Faden geschieht), dass die Distanz nicht grösser werden kann, während nichts
sie hindert, dass sie kleiner werde.

Es liessen sich noch andere Beispiele aufführen, wobei ich mich aber
nicht aufhalten will.

Die gewöhnliche Einkleidung schliesst nun, für sich betrachtet, alle solche
\clearpage\SourceLeaf{24}\HistoricalHead{19}{R}{GAUSS AN MÖBIUS.}
\noindent in der Natur so häufigen Fälle aus; \textls[120]{meine} Einkleidung umfasst sie auf einmal
alle mit.

Ich muss mich heute auf diese wenigen \textls[120]{Andeutungen} beschränken, aber
\textls[120]{sapienti sat}.

Es liessen sich noch mehrere andere Umstände, betreffend die ersten
Gründe der Statik und Bewegungslehre namhaft machen, die mir bei den
üblichen Behandlungen nicht tief genug aufgefasst zu sein scheinen. Allerdings
muss man dabei aber immer bedenken, für welchen Zweck und für
welche Leser geschrieben ist, und dass wo die Leser, für welche man schreibt,
keinen Anstoss nehmen, es vielleicht gar nicht wohlgethan wäre, tiefer einzudringen,
als ihnen frommt.

Aber für diejenigen, die gern Alles aus dem höchsten und allgemeinsten
Gesichtspunkt betrachten, scheint mir in allen Schriften zu wenig auf die
Unterscheidung der absoluten und der relativen Räume eingegangen zu werden.
Man hält sich von vornherein immer nur an die Phänomene auf unserer als
ruhend betrachtet werdenden kleinen Erde, ohne sich darüber zu rechtfertigen.
Ich gebe übrigens gern zu, dass es nicht ganz leicht ist, \textls[120]{gleich von vornherein}
Alles aus dem höheren Standpunkt eines über der materiellen Welt
stehenden Geistes darzustellen, und dass die Hülfe, welche die Haarspaltereien
der sogenannten Metaphysiker geben, nur eine sehr ungenügende sein möchte.

Sehen Sie alles Dieses nur so an, dass ich eine mit Mühe herausgerissene
halbe Stunde habe verwenden wollen, mich mit Ihnen über einige Ideenanregungen
etwas zu besprechen, wofür so wenige Menschen nur den Sinn
haben.

Unser Septembertermin wird morgen oder übermorgen wie gewöhnlich
abgewartet werden. Aber der nächste Novembertermin soll, da ich Herrn
\textsc{von Humboldt} damit einen Gefallen zu thun glaube, etwas versetzt und etwas
früher gehalten werden, etwa vom 13.—14. November oder 14.—15. November,
um zu sehen, ob in den berüchtigten Sternschnuppenzeiten bedeutende Magnetische
Bewegungen vorkommen. Die \textls[120]{genaue} Bestimmung des Tages (übrigens
nach gewöhnlicher Art 0\textsuperscript{h} bis 24\textsuperscript{h} Göttinger M. Z.) habe ich Herrn
\textsc{von Humboldt} \textls[120]{überlassen}, der sie nächstens durch die \textls[120]{Zeitungen} publiciren
wird. \textls[120]{In omnem eventum} wird es auch noch besonders brieflich den
Cooperatoren angezeigt werden.
\HistoricalFoot{}{3\textsuperscript{*}}

\clearpage\SourceLeaf{25}\HistoricalHead{20}{L}{PHYSIK. BRIEFWECHSEL.}
Erhalten Sie Ihr freundliches Andenken
\begin{flushright}Ihrem ergebensten\\C. F. \textsc{Gauss}.\end{flushright}
Göttingen den 29. September 1837.
\par\clearpage\endgroup
